\chapter{Scars}

\section*{Definition}
A scar is the macroscopic, fibrous tissue that replaces normal skin after injury or surgical incision, reflecting the end result of the mammalian wound healing continuum. Scar classification includes mature, immature, atrophic, hypertrophic, and keloid subtypes.

\section*{Pathophysiology}
Scar formation proceeds through three overlapping phases: inflammatory (hemostasis, neutrophil/macrophage infiltration), proliferative (fibroblast proliferation, collagen III deposition, angiogenesis, and granulation tissue), and remodeling (collagen III to collagen I cross-linking, myofibroblast-mediated contraction). In humans, scarred tissue lacks appendages (hair follicles, sweat glands), elastin fibers, and organized basket-weave collagen architecture, resulting in reduced tensile strength (max 80\% of intact skin) and a different texture/appearance from surrounding skin.

\section*{Causes}
\begin{itemize}
  \item \textbf{Cutaneous injury:} Surgical incisions, lacerations, burns, abrasions, puncture wounds
  \item \textbf{Inflammatory:} Acne vulgaris, chickenpox, folliculitis, hidradenitis suppurativa, vasculitic ulcers
  \item \textbf{Traumatic:} Blunt force injury (contusion/hematoma followed by fibrosis), abrasions, split-thickness skin graft donor sites
  \item \textbf{Surgical:} Elective procedures (laparotomy, breast surgery, joint replacement), laparoscopic port sites
  \item \textbf{Pathologic:} Hypertrophic scars (remain within wound boundaries, may regress), keloids (extend beyond margins, progressive), atrophic scars (depressed, as in acne or striae)
\end{itemize}

\section*{When to See a Doctor}
See your doctor if:
\begin{itemize}
  \item Scar is painful, pruritic, or hypertrophic beyond 3 months after injury
  \item Functional limitation: contracture across a joint restricting range of motion
  \item Cosmetic disfigurement causing psychosocial distress
  \item Ulceration or change in appearance of an old scar (rule out Marjolin ulcer --- squamous cell carcinoma arising in chronic scar)
  \item Keloid formation with progressive growth beyond wound margins
\end{itemize}

\section*{Related Symptoms}
\begin{itemize}
  \item Hypertrophic scar, keloid, contracture, wound dehiscence, atrophic scar, striae distensae
\end{itemize}
\textbf{Important:} The best treatment for an undesirable scar is prevention through meticulous wound closure, reduction of wound tension, and early silicone gel therapy beginning as soon as re-epithelialization is complete.

\section*{Self-Care / Home Management}
\begin{itemize}
  \item Apply silicone gel sheeting or silicone gel daily for at least 12 hours for 3--6 months on new scars
  \item Massage scar with emollient for 5--10 minutes daily to soften and flatten
  \item Sun protection (SPF 50+) on new scars for at least 6--12 months to prevent hyperpigmentation
\end{itemize}

\section*{How Your Doctor Will Evaluate You}
\begin{itemize}
  \item Clinical assessment: Scar type, color, contour, texture, size, pliability, pain, and pruritus
  \item Scar scales: Vancouver Scar Scale (VSS), POSAS (Patient and Observer Scar Assessment Scale) for objective measurement
  \item Biopsy if concerning features: ulceration, rapid growth, irregular borders, or suspected malignant transformation
\end{itemize}

\section*{Treatment Options}
\begin{itemize}
  \item Prevention: Tension-reducing closure techniques, silicone gel sheeting, pressure garments for burn scars
  \item Mature scars: Microneedling, laser therapy (pulsed-dye for erythema, CO2/fractional for texture), chemical peels
  \item Hypertrophic: Intralesional corticosteroid injection (triamcinolone 10--40 mg/mL) every 4--6 weeks, 5-FU combination
  \item Atrophic/acne scars: Subcision, microneedling, dermal fillers, TCA cross, fractional laser resurfacing
  \item Surgical revision: Z-plasty, W-plasty, geometric line closure, or scar excision with layered closure for unfavorable scars
  \item Radiation: Superficial radiotherapy for resistant keloids post-excision
\end{itemize}

\clearpage
